\documentclass[10pt,a4paper]{article}
\usepackage[left=2cm,right=2cm,top=2cm,bottom=2cm]{geometry}
\usepackage[T1]{fontenc}
\usepackage{lmodern}
\usepackage[spanish,es-tabla,es-noquoting]{babel}
\usepackage[dvipsnames]{xcolor}
\usepackage[fleqn]{mathtools}
\usepackage{booktabs}
\usepackage{amsmath}
\usepackage{latexsym}
\usepackage{graphicx}
\usepackage{nccmath}
\usepackage{multicol}
\usepackage{listings}
\usepackage{tasks}
\usepackage{color}
\usepackage{float}
\usepackage{tikz}
\usepackage{csquotes}
\usepackage[backend=biber,style=numeric,sorting=none]{biblatex}
\addbibresource{referencias.bib}

\definecolor{colorIPN}{rgb}{0.5, 0.0,0.13}
\definecolor{colorESCOM}{rgb}{0.0, 0.5,1.0}
\graphicspath{ {imagenes/} }

\begin{document}
%#########################################################
\begin{titlepage}
	\centering
	{ \huge \bfseries \color{colorIPN}{Instituto Politécnico Nacional} \par}
	{ \Large \bfseries  \color{colorESCOM}{Escuela Superior de Cómputo} \par }
	\vspace{1cm}
	{\huge\Large \color{colorIPN}{Web Client and Backend Development Frameworks}.\par}
	\vspace{1.5cm}
	{\huge\Large  \color{colorESCOM}{Tarea 1: Arquitectura de cebolla (arquitectura limpia)}\par}
		\vspace{2cm}
	{\Large\itshape \color{colorIPN}{Profesor: Dr. José Asunción Enríquez Zárate}\par} \hfill \break
	\vspace{2cm}
	{\Large\itshape \color{colorIPN}{Alumno: Angel Frausto Robles}\par} \hfill \break
	{\Large\itshape \color{colorIPN}{id634296@gmail.com}\par} \hfill \break
	{\Large\itshape \color{colorIPN}{7CM2} \par}
	\vfill
	{\large \color{colorIPN}{\today}\par}
	\vfill
\end{titlepage}

\renewcommand\lstlistingname{Código}


\lstset{
	language=Java,
	basicstyle=\small\sffamily,
	numbers=left,
	numberstyle=\tiny,
	frame=tb,
	tabsize=4,
	columns=fixed,
	showstringspaces=false,
	showtabs=false,
	keepspaces,
	commentstyle=\color{Violet},
	keywordstyle=\color{colorIPN} \bfseries,
	stringstyle=\color{colorESCOM},
	extendedchars=true,
	inputencoding=utf8,
	literate=
		{á}{{\'a}}1 {é}{{\'e}}1 {í}{{\'i}}1 {ó}{{\'o}}1 {ú}{{\'u}}1
		{Á}{{\'A}}1 {É}{{\'E}}1 {Í}{{\'I}}1 {Ó}{{\'O}}1 {Ú}{{\'U}}1
		{ñ}{{\~n}}1 {Ñ}{{\~N}}1 {ü}{{\"u}}1 {¿}{{?`}}1 {¡}{{!`}}1
}

\settasks{
	counter-format=(tsk[r]),
	label-width=4ex
}
\tableofcontents
\pagebreak
\listoffigures
\pagebreak
\listoftables
\pagenumbering {arabic}

\pagebreak

%################################################
\section{\color{colorIPN}{Introducción}}

El apartado 1.4 del temario corresponde a la arquitectura de cebolla, llamada también arquitectura limpia. Esta investigación cubre su origen, la regla que la define y los cuatro elementos que el propio temario enumera: modelo del dominio, interfaz de usuario, infraestructura y pruebas.

La arquitectura de cebolla responde a un problema concreto de la arquitectura tradicional de n capas, que el apartado 1.3 ya introduce. Ese problema es que la lógica de negocio acaba dependiendo de la base de datos. Lo que hace la cebolla es invertir esa dependencia, y de esa inversión se derivan tanto sus ventajas como su costo.

%################################################
\section{\color{colorIPN}{Conceptos previos}}

\subsection{ \color{colorESCOM}{El acoplamiento hacia la base de datos}}

En el esquema clásico de tres capas la interfaz de usuario llama a la capa de lógica de negocio, y ésta llama a la capa de acceso a datos. Cada capa depende de la que tiene debajo. Por transitividad, la lógica de negocio depende de la persistencia y la aplicación entera queda construida encima del esquema relacional.

Palermo describió la consecuencia en 2008 \cite{palermo2008}: cambiar el motor de almacenamiento, el ORM o el nombre de una columna obliga a tocar código de negocio que no tiene relación con esa decisión técnica. La base de datos, que es un detalle de almacenamiento, acaba funcionando como cimiento del sistema.

\subsection{ \color{colorESCOM}{El principio de inversión de dependencias}}

El principio de inversión de dependencias, la letra D de SOLID, enuncia que los módulos de alto nivel no deben depender de los de bajo nivel y que ambos deben depender de abstracciones. Aplicado a un par de clases, el principio produce una interfaz entre quien llama y quien ejecuta. Aplicado al sistema completo, produce la arquitectura de cebolla.

Dos nociones se confunden con frecuencia en este punto. La dependencia de código fuente es qué archivo necesita conocer a qué otro para compilar. El flujo de control es qué método llama a qué método durante la ejecución. La inversión consiste en hacer que ambos apunten en sentidos contrarios.

%################################################
\section{\color{colorIPN}{Desarrollo}}

\subsection{ \color{colorESCOM}{Origen y formulaciones}}

Jeffrey Palermo publicó en julio de 2008 la serie de artículos que dio nombre a la arquitectura de cebolla \cite{palermo2008}. Su propuesta se resume en cuatro postulados:

\begin{itemize}
	\item La aplicación se construye alrededor de un modelo de objetos independiente.
	\item Las capas internas definen interfaces y las externas las implementan.
	\item La dirección del acoplamiento apunta hacia el centro.
	\item Todo el código del núcleo puede compilarse y ejecutarse por separado de la infraestructura.
\end{itemize}

Cuatro años después, Robert C. Martin publicó el artículo que popularizó el nombre de arquitectura limpia \cite{martin2012} y en 2017 lo desarrolló en un libro \cite{martin2017}. Su enunciado central es la regla de dependencia: las dependencias del código fuente \enquote{\textit{can only point inwards}}. Ningún elemento de un círculo interior puede nombrar algo de un círculo exterior, y la restricción alcanza a clases, funciones, variables y formatos de datos.

Los dos autores describen la misma estructura con distinta nomenclatura. Palermo nombra sus anillos modelo del dominio, servicios de dominio, servicios de aplicación e infraestructura. Martin los nombra entidades, casos de uso, adaptadores de interfaz y frameworks. La figura \ref{img:cebolla} representa la versión de Palermo, que es la que sigue el temario.

\begin{figure}[H]
	\centering
	\begin{tikzpicture}
		\filldraw[fill=gray!5,  draw=colorESCOM, thick] (0,0) circle (4.1cm);
		\filldraw[fill=gray!12, draw=colorESCOM, thick] (0,0) circle (3.1cm);
		\filldraw[fill=gray!22, draw=colorESCOM, thick] (0,0) circle (2.1cm);
		\filldraw[fill=colorIPN!18, draw=colorIPN, very thick] (0,0) circle (1.15cm);

		\node[align=center, text=colorIPN] at (0,0)
			{\scriptsize\bfseries Modelo del\\[-2pt] \scriptsize\bfseries dominio};
		\node[fill=gray!22, inner sep=2pt] at (0,1.60)  {\scriptsize Servicios de dominio};
		\node[fill=gray!12, inner sep=2pt] at (0,2.60)  {\scriptsize Servicios de aplicación};
		\node[fill=gray!5,  inner sep=2pt] at (0,3.60)  {\scriptsize Interfaz de usuario};
		\node[fill=gray!5,  inner sep=2pt] at (-2.30,-2.77) {\scriptsize Infraestructura};
		\node[fill=gray!5,  inner sep=2pt] at (2.30,-2.77)  {\scriptsize Pruebas};

		\draw[->, very thick, colorIPN] (3.95,0) -- (1.35,0);
	\end{tikzpicture}
	\caption{Anillos de la arquitectura de cebolla y dirección de las dependencias}
	\label{img:cebolla}
\end{figure}

\subsection{ \color{colorESCOM}{Modelo del dominio}}

El modelo del dominio ocupa el centro. Contiene las entidades y los objetos de valor que representan las reglas del negocio, aquellas que seguirían siendo ciertas aunque la aplicación no existiera como software. Entran aquí un pedido que calcula su total a partir de sus líneas, o una cuenta que rechaza cualquier operación que la deje por debajo del saldo mínimo. El vocabulario y la técnica provienen del diseño guiado por el dominio de Evans \cite{evans2003}.

En Java esto significa clases sin anotaciones de framework, sin \texttt{@Entity}, sin \texttt{@Component} y sin importaciones de Spring ni de JPA. La restricción admite una comprobación objetiva: si el módulo del dominio compila con un descriptor de construcción que no declara ninguna dependencia de infraestructura, el núcleo está aislado \cite{foojay2023}, y si no compila existe una fuga que hay que localizar.

Alrededor del modelo, Palermo coloca los servicios de dominio, que resuelven operaciones de negocio que no pertenecen a una sola entidad, y el primer anillo de interfaces, donde viven los repositorios. Un repositorio se declara dentro del núcleo como interfaz y se implementa fuera de él, de modo que el núcleo declara qué necesita mientras la forma de satisfacerlo queda en el exterior. El código \ref{lst:puerto} muestra esa separación.

\begin{lstlisting}[caption={Interfaz declarada en el núcleo y caso de uso que la consume},label={lst:puerto}]
// Nucleo: la interfaz vive junto al dominio
public interface PedidoRepository {
	Optional<Pedido> buscarPorId(PedidoId id);
	void guardar(Pedido pedido);
}

// Caso de uso: depende de la abstraccion, no del ORM
public class ConfirmarPedido {

	private final PedidoRepository repositorio;

	public ConfirmarPedido(PedidoRepository repositorio) {
		this.repositorio = repositorio;
	}

	public void ejecutar(PedidoId id) {
		Pedido pedido = repositorio.buscarPorId(id)
				.orElseThrow(PedidoNoEncontrado::new);
		pedido.confirmar();
		repositorio.guardar(pedido);
	}
}
\end{lstlisting}

Ninguna de las dos clases menciona una tecnología de persistencia. La clase \texttt{ConfirmarPedido} corresponde al anillo de servicios de aplicación en la terminología de Palermo y al anillo de casos de uso en la de Martin. Orquesta entidades para cumplir un objetivo del usuario y decide la secuencia de pasos, sin saber de dónde llegan los datos ni dónde se guardan.

\subsection{ \color{colorESCOM}{Interfaz de usuario}}

La interfaz de usuario se sitúa en el anillo exterior, al mismo nivel que la infraestructura y las pruebas. Funciona como mecanismo de entrega y por eso es reemplazable: la misma aplicación puede exponerse como sitio web, como API consumida por un cliente de escritorio o como línea de comandos, sin que el núcleo cambie.

En la pila que usa el curso, un cliente Angular consume una API REST publicada por controladores de Spring, y ambas piezas viven en el anillo exterior. El controlador recibe la petición HTTP, la traduce a una llamada al caso de uso y convierte el resultado en JSON, sin que el dominio necesite saber que HTTP existe.

Martin llama a este anillo adaptadores de interfaz y le asigna los controladores, los presentadores y los conversores de formato \cite{martin2012}. De ahí sale un criterio práctico para ubicar una clase: si menciona \texttt{HttpServletRequest}, \texttt{ResponseEntity} o un DTO pensado para serializarse en JSON, pertenece al anillo exterior y no al núcleo.

La dirección del acoplamiento se conserva aunque sea el usuario quien inicia la interacción. El controlador conoce al caso de uso y lo llama, mientras que el caso de uso ignora que el controlador existe.

\subsection{ \color{colorESCOM}{Infraestructura}}

El anillo de infraestructura agrupa todo lo que comunica a la aplicación con el exterior: persistencia, mensajería, correo, clientes HTTP hacia servicios de terceros, sistema de archivos y relojes del sistema. Su trabajo consiste en implementar las interfaces que el núcleo declaró.

En tiempo de compilación el módulo de infraestructura depende del núcleo, porque necesita conocer la interfaz \texttt{PedidoRepository} para poder implementarla. En tiempo de ejecución el flujo de control va en sentido contrario, ya que el caso de uso invoca al repositorio y termina ejecutando código de JPA. Las dos flechas apuntan en direcciones opuestas, y ese cruce es todo el contenido del principio de inversión de dependencias.

El ensamblado ocurre en un solo punto, la raíz de composición. En Spring corresponde a una clase de configuración o al escaneo de componentes, y es el único lugar del sistema donde se nombran a la vez la abstracción y la implementación concreta que la satisface \cite{foojay2023}.

Sustituir PostgreSQL por MongoDB modifica el módulo de infraestructura y la configuración, y nada más. El dominio y los casos de uso no se recompilan por ese cambio, y sus pruebas siguen pasando sin tocarlas.

\subsection{ \color{colorESCOM}{Pruebas}}

El diagrama original de Palermo coloca las pruebas en el anillo exterior, junto a la interfaz de usuario y la infraestructura \cite{palermo2008}. Las sitúa ahí porque una prueba es otro cliente del núcleo, que ocupa el mismo lugar que la interfaz de usuario y entra al sistema por las mismas fronteras.

Como el núcleo no depende de frameworks, sus pruebas son unitarias en sentido estricto y no requieren base de datos, contexto de Spring ni servidor de aplicaciones. Se construye la entidad, se invoca el método y se verifica el invariante. La testabilidad figura entre los objetivos declarados de la arquitectura \cite{martin2012} porque se sigue directamente de la regla de dependencia.

La distribución de pruebas suele seguir una pirámide organizada por anillos:

\begin{itemize}
	\item Las pruebas del dominio son las más numerosas. Verifican invariantes, cálculos y transiciones de estado, y corren rápido porque no tocan nada externo.
	\item Cada caso de uso tiene sus propias pruebas. Como los repositorios son interfaces del núcleo, se sustituyen por implementaciones en memoria o por dobles.
	\item La infraestructura se prueba por integración contra la base de datos real, habitualmente con contenedores efímeros, para comprobar que el adaptador cumple el contrato que declaró.
	\item Las pruebas de extremo a extremo son las menos, por costo de mantenimiento y tiempo de ejecución.
\end{itemize}

Bogdanović documenta un error frecuente al aplicar este esquema \cite{optivem2023}: los dobles deben sustituir dependencias externas al núcleo, no al dominio mismo. Cuando se reemplazan entidades por mocks, las pruebas dejan de verificar comportamiento y pasan a verificar llamadas, de modo que cualquier refactorización interna las rompe aunque el sistema siga siendo correcto.

%################################################
\section{\color{colorIPN}{Análisis comparativo}}

\subsection{	\color{colorESCOM}{Tres nombres para la misma regla}}

La arquitectura de cebolla tuvo un antecedente directo. Alistair Cockburn había publicado en 2005 la arquitectura hexagonal, también llamada de puertos y adaptadores \cite{cockburn2005}, y el propio Palermo la reconoce como tal. La tabla \ref{tabla:comparacion} resume las tres formulaciones.

\begin{table}[H]
	\centering
	\begin{tabular}{p{2.4cm} |p{3.8cm} |p{3.8cm} |p{4.2cm}}\hline \hline
		& \textbf{Hexagonal (2005)} & \textbf{Cebolla (2008)} & \textbf{Limpia (2012)} \\ \hline \hline
		Autor & Alistair Cockburn & Jeffrey Palermo & Robert C. Martin \\ \hline
		Metáfora & Hexágono rodeado de puertos y adaptadores & Anillos concéntricos & Círculos concéntricos \\ \hline
		Centro & La aplicación & Modelo del dominio & Entidades \\ \hline
		Divisiones internas & No las prescribe & Dominio, servicios de dominio, servicios de aplicación & Entidades, casos de uso, adaptadores, frameworks \\ \hline
		Énfasis propio & Simetría entre el lado que entra y el que sale & Disciplina de anillos y servicios de dominio & Casos de uso explícitos y enunciado de la regla de dependencia \\ \hline \hline
	\end{tabular}
	\caption{Comparación entre las tres arquitecturas centradas en el dominio}
	\label{tabla:comparacion}
\end{table}

Las tres se apoyan en el principio de inversión de dependencias, ponen la lógica de negocio en el centro, dejan la tecnología en el borde y obtienen testabilidad como resultado \cite{graca2017}. Jovanović las comparó capa por capa y concluyó que la estructura de carpetas a la que llegan los tres enfoques es prácticamente la misma, y que la discusión sobre cuál adoptar importa menos que aplicar el criterio de forma consistente \cite{jovanovic2024}.

\subsection{	\color{colorESCOM}{Costos y críticas}}

El costo principal aparece en el mapeo entre fronteras. Un mismo concepto de negocio suele existir por triplicado: como entidad del dominio, como DTO de entrada o salida y como modelo de persistencia. Cada frontera exige un traductor, y ese código de traducción hay que escribirlo y mantenerlo.

Three Dots Labs recoge los señalamientos habituales \cite{threedots2023}, que coinciden entre sí: interfaces con una sola implementación que nunca tendrán otra, capas que se limitan a reenviar la llamada a la siguiente sin agregar nada, y código de una misma funcionalidad repartido entre carpetas separadas por criterio técnico, de modo que entender una operación completa obliga a recorrer varios módulos. A esto se suma la indirección adicional en tiempo de ejecución, que en sistemas sensibles al rendimiento hay que medir en lugar de suponer.

El criterio de decisión pasa por la complejidad de las reglas de negocio y por el tiempo de vida esperado del sistema, más que por el tamaño del código. Una aplicación que traslada formularios a tablas sin reglas propias no gana nada con cuatro anillos. Un sistema cuyas reglas cambian, con varios equipos trabajando en paralelo y años de mantenimiento por delante, amortiza el costo de la separación.

%################################################
\section{\color{colorIPN}{Conclusión}}

La arquitectura de cebolla aplica el principio de inversión de dependencias a la escala del sistema completo. En lugar de invertir la relación entre dos clases, invierte la relación entre la lógica de negocio y la tecnología que la rodea, de manera que la base de datos, el framework web y el cliente pasan a ser detalles sustituibles en el borde.

El resultado admite comprobación objetiva. El núcleo compila sin infraestructura en el classpath, sus pruebas corren sin base de datos, y cambiar de motor de almacenamiento afecta a un módulo identificable. Cuando alguna de esas afirmaciones deja de cumplirse, la regla de dependencia se rompió en algún punto del código.

Las formulaciones hexagonal, de cebolla y limpia difieren en nomenclatura y en el grado de detalle con que prescriben los anillos interiores, pero coinciden en la regla y en la estructura resultante. Elegir entre ellas importa menos que sostener la dirección de las dependencias de forma consistente durante todo el proyecto.

Para el proyecto del curso la traducción es directa. Las reglas de negocio se escriben en Java sin anotaciones, Spring y JPA quedan confinados al anillo exterior, Angular es un mecanismo de entrega más, y una raíz de composición se encarga de unir las piezas.


%################################################
\pagebreak

\section{\color{colorIPN}{Referencias Bibliográficas}}
\color{colorESCOM}{
	\printbibliography[heading=none]
}

\end{document}
